%-------------------------
% Cover Letter in Latex
% Enrique Martínez Martel
%------------------------

\documentclass[letterpaper,11pt]{article}

\usepackage{latexsym}
\usepackage[empty]{fullpage}
\usepackage{marvosym}
\usepackage[usenames,dvipsnames]{color}
\usepackage[hidelinks]{hyperref}
\usepackage{fancyhdr}
\usepackage[english]{babel}
\usepackage{tabularx}
\usepackage{ragged2e}
\usepackage{fontawesome5}
\usepackage{graphicx}
\usepackage[scale=0.90,lf]{FiraMono}
\usepackage{enumitem}

\definecolor{light-grey}{gray}{0.83}
\definecolor{dark-grey}{gray}{0.3}
\definecolor{text-grey}{gray}{.08}

\usepackage{contour}
\usepackage[normalem]{ulem}
\renewcommand{\ULdepth}{1.8pt}
\contourlength{0.8pt}
\newcommand{\myuline}[1]{%
  \uline{\phantom{#1}}%
  \llap{\contour{white}{#1}}%
}

\usepackage{tgheros}
\renewcommand*\familydefault{\sfdefault}
\usepackage[T1]{fontenc}

\pagestyle{fancy}
\fancyhf{}
\fancyfoot{}
\renewcommand{\headrulewidth}{0pt}
\renewcommand{\footrulewidth}{0pt}

\addtolength{\oddsidemargin}{-0.5in}
\addtolength{\evensidemargin}{0in}
\addtolength{\textwidth}{1in}
\addtolength{\topmargin}{-.8in}
\addtolength{\textheight}{1.5in}

\urlstyle{same}
\raggedbottom
\raggedright

\color{text-grey}

\begin{document}

\vspace*{160pt}

Dear LTS4 selection committee,

\vspace{16pt}

I came across the SATURNA initiative and it stood out immediately. Not many places are combining AI with RNA biology at this level of ambition, and the LTS4 lab at EPFL is one of the best environments in the world to do it. That is what brought me here. My background is in deep learning with a focus on transformer architectures and large language models: a multimodal LLM classification pipeline built at Banco Santander evaluating GPT-4/GPT-4o on real operational data, an MSc thesis at UPM benchmarking deep learning architectures on a real-world dataset (grade: 10/10), and a BSc thesis on autonomous perception at UPC. More of my work is at \href{https://enkairito.github.io/enriquemm.github.io/}{\myuline{enkairito.github.io/enriquemm.github.io}}.

\vspace{16pt}

The main reasons I want to do this PhD:

\vspace{8pt}

\begin{itemize}[leftmargin=1.5em, itemsep=10pt]

  \item \textbf{Language models applied to sequences that encode life.} RNA is a sequence. I know how to model sequences. But the interesting part is not the overlap. What I find exciting about SATURNA is the ambition to build models that understand why a sequence matters: what it does, how it breaks, and what happens when it does. Generic language models were not built to answer those questions.

  \item \textbf{AI co-developed with domain experts.} The idea of working side by side with biologists, neuroscientists, and clinicians, people who see the problem from a completely different angle, is something that genuinely excites me. I have always found that the best work happens when different kinds of expertise are in the same room, and SATURNA is built exactly around that.

  \item \textbf{Clinical translation as ground truth.} In industry you spend a lot of time hitting targets that are one step removed from what anyone really cares about. Here the question is whether a patient gets better. That is something I genuinely want to contribute to, using what I know about AI to help people who are actually going through something difficult.

\end{itemize}

\vspace{16pt}

I do not have a biology or RNA background. I am not pretending otherwise. But sequence modeling, generative architectures, and learning from structured biological data are areas I am actively developing, and I learn new domains fast when the problem is worth understanding.

\vspace{16pt}

I would be happy to discuss further.

\vspace{20pt}

Best regards,

\vspace{16pt}

\textbf{Enrique Martínez Martel}

\vspace{6pt}

\faEnvelope \hspace{2pt} \texttt{enri.martinez25@gmail.com} \hspace{6pt} $|$ \hspace{6pt}
\faLinkedin \hspace{2pt} \href{https://www.linkedin.com/in/enrique-martinez-martel}{\myuline{enrique-martinez-martel}} \hspace{6pt} $|$ \hspace{6pt}
\faGlobe \hspace{2pt} \href{https://enkairito.github.io/enriquemm.github.io/}{\myuline{enkairito.github.io/enriquemm.github.io}}

\end{document}
